% Generated by `taskgen course examples` from dlp.rho@sem2 (seed 21).
% Do not edit: change examples.yaml and regenerate.
\begin{example}[$\rho$-алгоритм Полларда с циклом Флойда]\label{ex:rho-floyd}
Решить задачу дискретного логарифмирования $114^x \equiv 372\ (\textrm{mod}\ 569)$ $\rho$-алгоритмом Полларда с поиском цикла по Флойду; $\ord(114) = 71$. Последовательность $x_{i+1} = f(x_i)$, $x_i = 114^{\alpha_i} \cdot 372^{\beta_i} \bmod 569$ (показатели по модулю $71$), где\[f(x) = \begin{cases} x^2, & x \equiv 0,\\ 372\, x, & x \equiv 1,\\ 114\, x, & x \equiv 2 \end{cases} \pmod 3, \qquad (\alpha, \beta) \mapsto \begin{cases} (2\alpha, 2\beta),\\ (\alpha, \beta + 1),\\ (\alpha + 1, \beta), \end{cases}\]начиная с $x_0 = 114^{10} \cdot 372^{6} \equiv 1$ ($\alpha_0 = 10$, $\beta_0 = 6$). Выписать таблицу $(x_i, \alpha_i, \beta_i)$, $(x_{2i}, \alpha_{2i}, \beta_{2i})$ до совпадения $x_i = x_{2i}$ и найти $x$.

\par\noindent\textit{Решение.} Вычисляем «черепаху» $x_i = f(x_{i-1})$ и «зайца» $x_{2i} = f(f(x_{2i-2}))$:\[\begin{array}{r|rrr|rrr}i & x_i & \alpha_i & \beta_i & x_{2i} & \alpha_{2i} & \beta_{2i} \\ \hline 0 & 1 & 10 & 6 & 1 & 10 & 6 \\ 1 & 372 & 10 & 7 & 117 & 20 & 14 \\ 2 & 117 & 20 & 14 & 520 & 9 & 56 \\ 3 & 33 & 40 & 28 & 400 & 18 & 43 \\ 4 & 520 & 9 & 56 & 469 & 36 & 17 \\ 5 & 549 & 9 & 57 & 136 & 1 & 36 \\ 6 & 400 & 18 & 43 & 549 & 1 & 38 \\ 7 & 291 & 18 & 44 & 291 & 2 & 6 \\ \end{array}\]Совпадение при $i = 7$: $x_{7} = x_{14} = 291$ (у орбиты $x_0$ предпериод $T = 4$ и период $M = 7$; $i = 7$ --- наименьшее кратное $M$, не меньшее $T$).\par Из $g^{18} h^{44} = g^{2} h^{6}$ и $h = g^x$ получаем $(44 - 6)\, x \equiv 2 - 18\ (\textrm{mod}\ 71)$, т.\,е. $38\, x \equiv 55\ (\textrm{mod}\ 71)$. $\gcd(38, 71) = 1$, решение единственно: $x \equiv 55 \cdot 38^{-1}$.  Обратный к $38$ по модулю $71$ (расширенный алгоритм Евклида для $(71, 38)$):
\[\begin{array}{r|r|r|r}
r_i & q_i & u_i & v_i \\ \hline
71 &  & 1 & 0 \\
38 &  & 0 & 1 \\
33 & 1 & 1 & -1 \\
5 & 1 & -1 & 2 \\
3 & 6 & 7 & -13 \\
2 & 1 & -8 & 15 \\
1 & 1 & 15 & -28 \\
0 & 2 & -38 & 71 \\
\end{array}\qquad r_i = u_i \cdot 71 + v_i \cdot 38\]

$38^{-1} \equiv -28 \equiv 43\ (\textrm{mod}\ 71)$.
\par $\tilde x \equiv 55 \cdot 43 \equiv 22\ (\textrm{mod}\ 71)$. Проверка: $114^{22} \equiv 372\ (\textrm{mod}\ 569)$.\[x = 22.\]
\end{example}
